\section{SFT Dataset Design: además de la cantidad, la distribución}

La sección anterior eligió un human-written baseline; esta sección aborda el problema de adquisición realmente difícil: cuántos ejemplos se necesitan, de qué tareas, de qué longitud y qué población los redacta. La aportación central de la clase es reformular «recolectar 10K ejemplos» como un design vector multidimensional.

\subsection{Diminishing Returns no significa «los datos no importan»}

La sección anterior explicó por qué H4 necesita un human-written baseline; esta subsección pregunta además, dado que los datos escritos por personas son caros, en qué punto aumentar el número de ejemplos empieza a perder valor marginal. Al leer la curva, primero hay que distinguir entre «la pendiente del beneficio disminuye» y «los datos dejan de importar». Trabajos previos sugieren que el beneficio de las instruction de alta calidad se desacelera después de unos cuantos miles de ejemplos, pero esta conclusión depende del modelo base, de la mezcla de tareas y del evaluator. La interpretación correcta es que el beneficio marginal decrece y que hace falta una controlled ablation; la interpretación errónea es que cualquier conjunto de 1K ejemplos basta.

\lecturefigure{slide-14.jpg}{Escala de datos y diminishing returns en investigaciones previas sobre SFT}{Diapositiva oficial p.14}

\begin{knowledgebox}{Lectura de la figura: la curva respalda un beneficio marginal decreciente, no un tamaño óptimo universal}
El eje horizontal es el número de instruction añadidas y el eje vertical es el desempeño en una evaluación específica. Primero observa en qué momento disminuye la pendiente y después revisa el task mix y el model que usa la curva. La saturación indica que seguir replicando ejemplos similares aporta menos valor, pero no descarta que añadir tareas nuevas, aumentar la dificultad o corregir errores siga siendo eficaz.
\end{knowledgebox}

\teachervoice{Al iniciar la adquisición, el equipo partía de dos supuestos previos aparentemente contradictorios: la escala de datos debería estar en tens of thousands, pero después de unas pocas instructions de alta calidad el desempeño se satura rápidamente. El ponente usa esta contradicción para introducir las task/length ablation posteriores, en lugar de ofrecer un número mágico de ejemplos.}

\lecturefigure{slide-15.jpg}{SFT data desiderata: task, length, quantity, human y synthetic}{Diapositiva oficial p.15}

\begin{importantbox}{Dataset size no es un estadístico suficiente}
Dos conjuntos de datos con 10K ejemplos cada uno pueden diferir por completo en task entropy, prompt length, completion length, multi-turn ratio, perfil de los redactores y correlación de los errores. Si los modelos muestran desempeños distintos, la causa no puede atribuirse únicamente a ``more data'' o ``better quality''.
\end{importantbox}

\subsection{Task Distribution: qué practica finalmente el modelo}

El task mix determina cómo se reparte el token budget entre capacidades como generation, QA, coding y classification. La clase muestra primero la distribución de InstructGPT y después la que H4 eligió para la Surge collection, y subraya que la distribution es una función conjunta de los requisitos del producto y de la posibilidad de evaluarlo.

\lecturefigure{slide-16.jpg}{Instruction task distribution de InstructGPT}{Diapositiva oficial p.16}

\begin{knowledgebox}{Lectura de la figura: los porcentajes son una asignación del presupuesto de entrenamiento}
Cada fila de la tabla no solo representa el nombre de una tarea, sino también un espacio de respuestas distinto: classification tiene entropía más baja, brainstorm/generation tiene entropía más alta, y en coding/math es más probable que exista una corrección ejecutable o única. Un cambio en la distribución modifica a la vez la dificultad de aprendizaje y la evaluator reliability.
\end{knowledgebox}

\lecturefigure{slide-17.jpg}{Task mix elegido por H4/Surge}{Diapositiva oficial p.17}

\begin{knowledgebox}{Lectura de la figura: comparar diferencias, no buscar la única receta correcta}
Primero compara cómo se desplazan las proporciones de generation, rewrite, coding, QA, etc., respecto a la línea base, y después pregunta a qué objetivo sirve ese desplazamiento. La distribución de H4 es un diseño con un presupuesto, una población de redactores y unos supuestos de aplicación de 2023; no es la universal mixture de todos los chatbot.
\end{knowledgebox}

\teachervoice{El ponente resume el problema de adquisición como «how many thousands of what task?». Esta frase es más precisa que «la calidad de los datos importa más»: la calidad tiene que concretarse en cobertura de tareas, posibilidad de determinar si una respuesta es correcta y objetivos de las tareas posteriores.}

\begin{lstlisting}[caption={Código conceptual para repartir un demonstration budget limitado según proporciones objetivo}]
target_mix = {
    "generation": 0.15, "open_qa": 0.05, "brainstorm": 0.10,
    "rewrite": 0.15, "summarize": 0.10, "math": 0.05,
    "coding": 0.15, "classify": 0.10, "closed_qa": 0.05,
    "extract": 0.10,
}

def allocate_examples(total_examples, target_mix):
    counts = {task: round(total_examples * ratio)
              for task, ratio in target_mix.items()}
    counts[max(counts, key=counts.get)] += total_examples - sum(counts.values())
    return counts
\end{lstlisting}

\subsection{Length Distribution: más largo no significa más útil}

La prompt/response length modifica la task difficulty, el número de tokens de entrenamiento, el generation style y las preferencias del judge. La clase coloca lado a lado la length table de los prior datasets y las restricciones elegidas, lo que anticipa la metric reversal que aparece más adelante.

\lecturefigure{slide-18.jpg}{Distribución de prompt y completion length en conjuntos de datos existentes}{Diapositiva oficial p.18}

\begin{knowledgebox}{Lectura de la figura: el número de ejemplos y el número de tokens son dos presupuestos distintos}
Con los mismos 10K examples, duplicar la longitud media de completion consume aproximadamente el doble de target tokens. Las respuestas largas también tienden a incluir listas, explicaciones y redundancia, y un LLM judge puede tomar estas style signal como helpfulness; por eso length es a la vez una variable de entrenamiento y un evaluation confounder.
\end{knowledgebox}

\lecturefigure{slide-19.jpg}{Length constraint elegida por H4 e intervalos prioritarios}{Diapositiva oficial p.19}

\begin{warningbox}{El length effect no puede juzgarse solo con curvas de correlación}
Si los prompt largos se concentran en ciertas tareas, o las completion largas provienen de redactores específicos, la performance--length correlation puede deberse a task/source confounding. El valor de las ablation posteriores consiste en controlar parte de las variables, pero aun así no permite obtener automáticamente una ley causal universal.
\end{warningbox}

\subsection{Tareas, ejemplos y población de redactores de Surge-Instruct}

Las dos subsecciones anteriores ya fijaron dos variables de diseño, task y length; esta subsección las aplica al conjunto de datos finalmente adquirido y añade la workforce provenance, que suele pasarse por alto. Al leer las tres diapositivas siguientes, responde en orden «qué se recolectó», «cómo son los ejemplos» y «quién los escribe», y después evalúa cómo esta información limita la extrapolación de los resultados. Surge-Instruct recolectó finalmente unos 10K instruction--demonstration pairs; la descripción del conjunto de datos solo está completa cuando se reúnen los tres tipos de evidencia.

\lecturefigure{slide-20.jpg}{Composición por tareas de los 10K datos de Surge-Instruct}{Diapositiva oficial p.20}

\begin{knowledgebox}{Lectura de la figura: el gráfico de pastel y la tabla de conteos deben leerse juntos}
Generation ocupa la porción más grande, pero open QA, brainstorm, chat, rewrite, summarize, coding, etc., forman en conjunto la mixture. Primero revisa los conteos absolutos para que el gráfico de pastel no oculte las categorías pequeñas, y luego verifica si las tareas cubren el producto previsto; la figura describe la distribution, pero no demuestra directamente qué categoría produce la mejora.
\end{knowledgebox}

\lecturefigure{slide-21.jpg}{Ejemplos de generation, classification, brainstorm y QA en Surge-Instruct}{Diapositiva oficial p.21}

\begin{knowledgebox}{Lectura de la figura: el criterio de «buena respuesta» cambia según la task}
Un knock-knock joke requiere estilo y estructura; classification requiere etiquetas discretas; brainstorm prioriza la cobertura y la concisión; open QA prioriza los hechos. Entrenarlos juntos con la misma loss es sencillo, pero al evaluar hay que conservar task-specific criteria.
\end{knowledgebox}

\lecturefigure{slide-22.jpg}{Perfil demográfico de la Surge annotator workforce}{Diapositiva oficial p.22}

\begin{warningbox}{Demographics es una descripción de la procedencia, no una prueba de representatividad}
La US-based workforce y su distribución por género, edad, raza y escolaridad ayudan a entender quién produjo los datos, pero no demuestran que representen a los usuarios de todo el mundo. Las notas conservan esta información para dar seguimiento a la viewpoint coverage y a los posibles sesgos, no para tratar la demographic diversity como una certificación de calidad.
\end{warningbox}

\teachervoice{Al enumerar el annotator background, el ponente reconoce que la human-written data también tiene una distribución de producción. Los datos escritos por personas no carecen de sesgo de manera automática; solo trasladan la fuente del sesgo del teacher model a los redactores, las guías, el proveedor y el proceso de revisión.}

\subsection{Resumen de la sección}

La SFT collection es un diseño conjunto de task, length, quantity y workforce. Solo si se registran estas variables pueden explicarse las diferencias posteriores en el leaderboard; informar únicamente «10K datos de alta calidad» descarta la información experimental más importante.
